\chapter*{How to Use This Document}
\addcontentsline{toc}{chapter}{How to Use This Document}

This reference document accompanies the EDNS~491 and EDNS~492 course syllabi and the Canvas course site. It is organized to support you across the entire two-semester capstone design experience, from your first client meeting through the Design Showcase poster session.

\textbf{Three documents work together:}
\begin{enumerate}[nosep,leftmargin=1.5em]
\item \textbf{This Course Text}: deep technical and process content organized in five parts. Consult it when you need the methodology, design procedures, project management frameworks, communication guidance, technical reference, or safety protocols. Start with the ``What Should I Be Reading Right Now?'' triage guide on the next page.
\item \textbf{The Course Syllabus}: assignment details, grading weights, policies, CLOs, and faculty contact information. The syllabus is the authoritative source for deadlines and grading.
\item \textbf{Canvas}: assignment rubrics, submission portals, scheduling, lecture materials, professional development modules, purchasing forms, and the Labriola InnoHub enrollment site.
\end{enumerate}

\textbf{Part~I: The Design Process} (Chapters~1--6) covers the complete engineering design lifecycle from problem definition through verification and validation. These chapters support your work in chronological order as your project progresses through design reviews.

\textbf{Part~II: Major Project Milestones} (Chapters~7--10) provides detailed guidance for the Statement of Work, PDR, CDR, and FDR, with deliverable checklists, evaluation criteria, and common failure patterns at each gate.

\textbf{Part~III: Project Management and Operations} (Chapters~11--16) covers hybrid project management, fabrication spaces and resources, team dynamics, budgeting, purchasing procedures, and travel planning. These chapters are practical operations guides that you will reference repeatedly throughout both semesters.

\textbf{Part~IV: Professional Skills} (Chapters~17--20) covers client engagement, written and oral communication, and career preparation. These skills distinguish capstone from prior coursework and directly affect your individual performance evaluations.

\textbf{Part~V: Safety, Compliance, and Professional Practice} (Chapters~21--22) covers EHS requirements, professional ethics, PE licensure, and course policies.

\textbf{Appendices} include a glossary, quick reference tables, ``Why Capstone Projects Fail'' (12 common patterns), ABET outcome mapping, templates and checklists, the NCEES Rules of Professional Conduct, and references.

\subsection*{Reading by Role}
\begin{itemize}[nosep,leftmargin=1.5em]
\item \textbf{Scrum Master / Project Manager}: Chapters~11 (PM), 13 (Team Dynamics), 7 (SOW), 8--10 (Milestones)
\item \textbf{Chief Financial Officer}: Chapters~14 (Budgeting), 15 (Purchasing), 16 (Travel)
\item \textbf{Technical Lead}: Chapters~3 (Requirements), 4 (Concept Gen), 6 (V\&V), 9 (CDR)
\item \textbf{Communication Lead}: Chapters~17 (Client), 18 (Written), 19 (Oral), 13 (Team)
\item \textbf{Fabrication Lead}: Chapter~12 (Fab Spaces), 21 (EHS), Front Matter (Safety)
\end{itemize}

\subsection*{Project Tracks}
This document supports all three project tracks. Where guidance differs by track, look for \textbf{Build Track}, \textbf{Paper Track}, and \textbf{Competition Track} labels. All tracks are held to equivalent standards of engineering rigor. The Paper Track is not a reduced-effort alternative; it substitutes physical fabrication with analytical depth, simulation fidelity, and documentation completeness. The Competition Track does not reduce requirements; it reorganizes their timing around an external schedule.


% ========== TRIAGE TABLE ==========
